\begin{exercise}[Etude d'un alliage à mémoire de forme : le Nitinol]
%X-ENS, 2007

Un alliage à mémoire de forme (ou AMF) présente la propriété singulière de pouvoir mémoriser une forme déterminée au préalable. Un AMF existe à basse température sous une structure martensite se transformant sous l'effet d'un réchauffement en une structure austénite et vice-versa. On observe ainsi une transformation martensitique se produisant sans diffusion d'atomes, comme le montre la Fig.\ref{fig:Allo-CP-Ex03-Fig01} ci-dessous : 

\begin{center}
\includegraphics[width=0.5\textwidth]{/home/jeff/Documents/IUTOrsay/1-Cours/4p13-Alliages/Cours-TD/Exos/CP/Allo-CP-Ex03/figures/Fig01}
\captionof{figure}{Transition polymorphique \ce{austenite <=> martensite}}
\label{fig:Allo-CP-Ex03-Fig01}
\end{center}


Le nitinol est un AMF de nickel et de titane de composition st{\oe}chiométrique \ce{NiTi}. La phase austénite possède une structure cubique de type \ce{CsCl}. La phase martensite quant à elle adopte une structure cristalline, localement très proche de la phase austénite, ce qui assure la réversibilité de la transformation. %Le problème s'intéresse aux propriétés thermodynamiques de l'alliage et non à sa structure cristallographique. 

Dans tout l'énoncé, les grandeurs relatives à la martensite et à l'austénite seront respectivement indicées \ce{M} et \ce{A} : par exemple, \ce{H_M} désigne l'enthalpie molaire de la martensite. 

Le nitinol, mis en une forme $F_A$ à température élevée est ensuite refroidi. La transformation martensitique a alors lieu et on peut donner au solide obtenu une nouvelle forme $F_M$.

%1. Mélange idéal sous contrainte nulle.

\begin{questions}
\question Identifier la phase la plus ordonnée.

\question Exprimer l'enthalpie libre molaire de l'austénite et de la martensite en fonction
de la température.

\question Représenter les deux fonctions $G = f (T)$ pour l'austénite et la martensite. À quelle condition s'agit-il de droites ? Cette condition sera supposée vérifiée par la suite.%, sachant que l'austénite est la phase la plus désordonnée

\question Comment peut-on déterminer la température d'équilibre $T^\star$ des deux phases. Préciser les phases stables selon le domaine de température.

\question Quel est le signe de l'enthalpie de la transformation $\Delta_{A \to M} H$ ?

\question Représenter l'entropie du système contenant initialement $n$ mole d'austénite en fonction de la température. \`{A} la température $T^\star$ , quelle est l'entropie du système en fonction de la fraction molaire $x_A$ ?
\end{questions}

On peut naturellement obtenir ensuite $F_A$ à partir de $F_M$ en appliquant une contrainte à température constante ou en augmentant la température à contrainte nulle (ou constante).

\end{exercise}